% ============================================================================
% DISCUSSION
% ============================================================================
\section{Discussion}

\subsection{What Current Evidence Supports}
\label{sec:evidence_supports}

The experimental results provide evidence for the following properties
of the architecture.

\emph{Gradient-free learning.}  The MarkovTrie learns online from raw
byte sequences through surprise-modulated plasticity, adapting its own
decay rate, learning rate, and classification parameters without
backpropagation or differentiable loss surfaces.  The Field provides a
collective selection mechanism that modulates learning across the
network without computing gradients.

\emph{Modality-agnostic processing.}  The same byte-level encoding and
Value representation processes text, images, code, and structured
logical data through identical pipelines.  No modality-specific
encoders, tokenizers, or preprocessing stages are required.

\emph{Hardware-native parallelism.}  The Boolean instruction set maps
directly to SIMD and GPU architectures, with Boolean operations executing
at fixed, data-independent latency.  The multi-substrate backend
demonstrates that the same computation can be dispatched to CPU, CUDA,
or Metal without algorithmic modification.  The current implementation
extends that claim to the bounded geometric lane: CPU uses native
assembly for the PGA product, CUDA executes the lane in \texttt{float64},
and Metal executes an ABI-preserving \texttt{float32} shader path.

\emph{Geometric alignment.}  PGA Context multivectors now participate in
MarkovTrie transition rescoring, multimodal sensory/action Procrustes
alignment, and episodic memory re-alignment.  This provides an explicit
continuous operator where the earlier architecture relied only on
finite-field phase and affinity overlap.

\emph{Emergent specialization.}  The combination of SimHash affinity
projection and Kademlia routing produces content-based clustering:
Values with similar token content are routed to the same MarkovTrie,
which then specializes in that content domain.  This specialization
emerges from the routing geometry without explicit assignment.

\emph{Collective attention.}  The Field's eigenmode detection and
asymmetric projection demonstrate that attention-like behavior---
selective amplification of relevant patterns and suppression of
noise---can arise from gossip-propagated statistics and simple
threshold-based clustering, without query-key-value matrices or
softmax normalization.

\subsection{Known Limitations}
\label{sec:limitations}

\emph{Retrieval selectivity.}  The affinity-based routing is
locality-sensitive but permissive: Values sharing even modest bit
overlap are routed to overlapping node sets.  More selective gating
mechanisms---such as multi-stage filtering or learned hash
refinement---may be needed for high-precision discriminative retrieval.

\emph{Composition depth.}  Information propagation through sequences of
boolean operations attenuates with depth.  Each operation can only
redistribute existing bits or set bits to zero or one; there is no
amplification mechanism analogous to residual connections.  Deep
multi-step reasoning chains are therefore limited by the rate at which
signal degrades under repeated Boolean transformation.  The geometric
lane mitigates this for transition scoring but does not by itself supply
a general mechanism for deep symbolic composition.

\emph{Geometric validation.}  The PGA and Procrustes paths are now
implemented natively, but their contribution still needs systematic
ablation.  In particular, the runtime should quantify when the geometric
rescoring term improves beam ranking, when Procrustes improves
cross-domain action selection, and how much Metal's \texttt{float32}
shader arithmetic diverges from the CPU/CUDA \texttt{float64} reference.

\emph{Scale comparison.}  The system has not been benchmarked against
large autoregressive models on standard leaderboards.  The experimental
results establish baseline performance for the architecture but do not
claim competitive accuracy at the scale of contemporary language or
vision models.

\emph{Convergence.}  The adaptive self-tuning mechanism converges
empirically in the reported experiments but lacks formal convergence
guarantees.  Characterizing the conditions under which
surprise-modulated plasticity and eigenmode-based attention reliably
produce stable learning is an open problem.

\emph{Eigenmode stability.}  The greedy clustering algorithm for mode
detection is $O(nm)$ and does not guarantee optimal partitions.  Mode
boundaries may fluctuate across epochs, potentially causing oscillation
in the field pressure applied to borderline nodes.

\subsection{Directions for Future Work}
\label{sec:future_work}

\paragraph{Ablation Studies.}
Varying the frame width (512, 1024, 2048 bytes) would characterize the
trade-off between representational capacity and computational cost.
Disabling the Field and measuring the effect on downstream task
performance would isolate the contribution of eigenmode-based attention.
Varying the number of SimHash projections and the majority vote
parameter $k$ would quantify the relationship between affinity precision
and routing quality.  Disabling the geometric continuation boost,
domain Procrustes alignment, and episodic re-alignment independently
would quantify the contribution of continuous geometry.

\paragraph{Scaling Experiments.}
Measuring classification accuracy and generation quality as a function
of the number of Kadabra nodes ($10^1$ to $10^4$) would establish the
empirical scaling behavior of distributed learning.  Measuring trie
depth and node count as a function of corpus size would characterize the
memory growth profile of the MarkovTrie.

\paragraph{Field Dynamics.}
Formal analysis of the eigenmode detection and projection loop---
particularly conditions for convergence, oscillation bounds, and the
relationship between the coupling threshold $\theta_C$ and mode
stability---would strengthen the theoretical foundations of the
attention mechanism.

\paragraph{Structured Retrieval.}
The current system retrieves via trie lookup and episodic $k$-NN.
Introducing content-addressable indexing over the affinity space (e.g.,
locality-sensitive hashing trees) would enable sub-linear retrieval
across the full Kadabra network.

\paragraph{Geometric Kernel Measurement.}
The native kernels should be benchmarked across batch sizes and
substrates.  In particular, the CPU assembly path should be compared
against the scalar reference on both ARM64 and AMD64, CUDA should be
measured for candidate-batch throughput, and Metal should be evaluated
for numerical drift introduced by shader-side \texttt{float32}
arithmetic.
